

\subsubsection{Class II: Equilibrium Consistency}
The traces that fall under Class I are capable of exactly representing Saint~Venant's solutions in a standard finite element context, but without further restriction they may still exhibit certain undesirable properties.
In particular, the stress resultants... TODO
\begin{itemize}
    \item Tangent symmetry
    \item Nodal loads
\end{itemize}

The conjugate stiffness is constructed from the material strain energy:
\begin{equation}
\mathcal{U}(x) \coloneq \tfrac{1}{2}\int_\Section \bm{\varepsilon}\cdot \mathbb{C}\, \bm{\varepsilon}\dA 
= \tfrac{1}{2}\bm{e}\cdot \bar{\mathbf{C}}_T(x)\,\bm{e},
\label{eq:energy-conjugacy-def}
\end{equation}

Define the energetic tangent \(\IesanEnergy\):
\begin{equation}
\IesanEnergy \coloneq \int_\Section \mathbf{S}_0^\top \mathbb{C}\, \mathbf{S}_0\dA
\label{eq:def-iesan-energy}
\end{equation}

\subparagraph{S2}
\begin{definition}[Class II]\label{def:class-ii}
A Class I trace \(\mathcal{T}\) is a \textbf{Class II} trace if the energy-conjugate resultants $\bar{\bm{s}}$ satisfy the same equilibrium as the classical resultants \eqref{eq:equilibrium-ode}:
\begin{equation}
\bar{\bm{s}}' = \mathbb{J}\,\bar{\bm{s}}.
\label{eq:sbar-equilibrium}
\end{equation}
\end{definition}

\subparagraph{C2}

\begin{corollary}
Prismatic $\mathbf{C}_T$ implies prismatic $\bar{\mathbf{C}}_T$.
\end{corollary}
\begin{proof}
When $\TraceEvolveEP = \IesanEvolveSP$, we have $\mathbf{M}(x) = \mathbf{1}$ identically, 
which lies in $O(\IesanEnergy)$.
\end{proof}

\begin{remark}
The converse does not hold. 
Prismatic $\bar{\mathbf{C}}_T$ with non-prismatic 
$\mathbf{C}_T$ requires $\TraceEvolveEP - \IesanEvolveSP \neq \mathbf{0}$ 
but $\IesanEnergy$-skew. This is an exceptional situation. 
\end{remark}

\begin{proposition}[S2. Equilibrium-Conjugate]
\label{prop:equilibrium-compatibility}
Let $\mathbf{C}_T$ and $\bar{\mathbf{C}}_T$ be prismatic (C2). 
The condition 
\eqref{eq:sbar-equilibrium} holds if and only if
\begin{equation}
[\bar{\mathbf{G}}(x), \mathbb{J}] \stackrel{\rm S2}{=} \mathbf{0}, \qquad 
\bar{\mathbf{G}} \coloneq \bar{\mathbf{C}}_T(x)\,\mathbf{C}_T^{-1}(x)
\qquad\forall x.
\label{eq:commutator-condition}
\end{equation}
\end{proposition}

\begin{proof}
With prismatic stiffnesses, $\bar{\bm{s}} = \bar{\mathbf{C}}\,\bm{e}$ and 
$\snm = \mathbf{C}_T\,\bm{e}$, so $\bar{\bm{s}} = \bar{\mathbf{G}}\,\snm$. 
From \eqref{eq:equilibrium-ode}:
\begin{equation*}
\bar{\bm{s}}' = \bar{\mathbf{G}}\,\snm' = \bar{\mathbf{G}}\,\mathbb{J}\,\snm 
= \bar{\mathbf{G}}\,\mathbb{J}\,\bar{\mathbf{G}}^{-1}\,\bar{\bm{s}}.
\end{equation*}
For this to equal $\mathbb{J}\,\bar{\bm{s}}$, we require 
$\bar{\mathbf{G}}\,\mathbb{J}\,\bar{\mathbf{G}}^{-1} = \mathbb{J}$, i.e., 
$[\bar{\mathbf{G}}, \mathbb{J}] = \mathbf{0}$.
\end{proof}

The commutator condition \eqref{eq:commutator-condition} constrains the 
structure of $\bar{\mathbf{G}}$.

\begin{lemma}[Structure of the Centralizer]
\label{lem:centralizer}
A matrix $\bar{\mathbf{G}}$ commutes with the equilibrium generator $\mathbb{J}$ 
if and only if it has the block form
\begin{equation}
\bar{\mathbf{G}} = \begin{bmatrix}
\bar{G}_{11} & \bar{\bm{g}}_{12}^\top & \bar{G}_{13} & 0 \\
\mathbf{0} & \bar{\mathbf{G}}_{22} & \mathbf{0} & \mathbf{0} \\
\bar{G}_{31} & \bar{\bm{g}}_{32}^\top & \bar{G}_{33} & 0 \\
\bar{\bm{g}}_{41} & \bar{\mathbf{G}}_{42} & \bar{\bm{g}}_{43} & \bar{\mathbf{G}}_{44}
\end{bmatrix}
\label{eq:K-structure}
\end{equation}
with $\bar{\mathbf{G}}_{44} = -\FrameBasis^\times \bar{\mathbf{G}}_{22}\,\FrameBasis^\times$.
\end{lemma}

\begin{proof}
The commutator $[\bar{\mathbf{G}}, \mathbb{J}]$ has nonzero contributions only 
from the $(4,2)$ block of $\mathbb{J}$:
\begin{align*}
[\bar{\mathbf{G}}\mathbb{J}]_{i2} &= -\bar{G}_{i4}\,\FrameBasis^\times, \\
[\mathbb{J}\bar{\mathbf{G}}]_{4j} &= -\FrameBasis^\times \bar{G}_{2j}.
\end{align*}
Setting $[\bar{\mathbf{G}}, \mathbb{J}] = \mathbf{0}$:
\begin{itemize}
\item Rows 1--3, column 2: $\bar{G}_{14} = 0$, $\bar{\mathbf{G}}_{24} = \mathbf{0}$, $\bar{G}_{34} = 0$.
\item Row 2, columns 1, 3: $\bar{\mathbf{G}}_{21} = \mathbf{0}$, $\bar{\mathbf{G}}_{23} = \mathbf{0}$.
\item Row 4, column 2: $\bar{\mathbf{G}}_{44}\,\FrameBasis^\times = \FrameBasis^\times \bar{\mathbf{G}}_{22}$.
\end{itemize}
The constraint $\bar{\mathbf{G}}_{44}\,\FrameBasis^\times = \FrameBasis^\times \bar{\mathbf{G}}_{22}$ 
is equivalent to $\bar{\mathbf{G}}_{44} =-\FrameBasis^\times \bar{\mathbf{G}}_{22}\,\FrameBasis^\times$ 
since $(\FrameBasis^\times)^{2} = -\oneS$.
\end{proof}

The set of all matrices of the form \eqref{eq:K-structure} comprise the equilibrium invariance group $G$:
\[
G \coloneq \bigl\{\,\mathbf{K} \in \mathrm{GL}(6) \;\big|\; [\mathbf{K},\mathbb{J}] = 0\,\bigr\}.
\]
Every \(\mathbf{K} \in G\) has the block form \eqref{eq:K-structure}. 
This further restricts the space of trace matrices \(\TraceDeDpRoot\) to a nonlinear sub-manifold, for which a clean structure like that obtained under P1 does not exist.
However, it is possible to conveniently parameterize a subset of these traces, by first identifying a reference trace \(\mathcal{T}_R\) in the class. 

\begin{proposition}\label{prop:class-ii-e}
Consider a reference plane-section trace \(\mathcal{T}^{R}\) with trace matrix \(\TraceDeDpRoot^R\). 
% \begin{equation*}
%   \TraceDeDpRoot^R =
%   \begin{bmatrix}
%     1          & \bm{r}_R^\top & 0 & {\TraceNF}_R^\top \\
%     \mathbf{0} & \mathbf{0}    & {\TraceFT}_R & {\TraceFF}_R \\
%     0          & \mathbf{0}^\top & 1 & {\TraceTF}_R^\top \\
%     \mathbf{0} & -\FrameBasis^\times & {\TraceBT}_R & {\TraceMF}_R
%   \end{bmatrix}.
% \end{equation*}
Then \(\TraceDeDpRoot = \mathbf{X}\,\TraceDeDpRoot^R\) has the same plane-section
pattern
\begin{equation*}
  \TraceDeDpRoot =
  \begin{bmatrix}
    1          & \CenterTrace^\top & 0 & \TraceNF^\top \\
    \mathbf{0} & \mathbf{0}              & \TraceFT & \TraceFF \\
    0          & \mathbf{0}^\top         & 1 & \TraceTF^\top \\
    \mathbf{0} & -\FrameBasis^\times     & \TraceBT & \TraceMF
  \end{bmatrix},
\end{equation*}
with blocks expressed in terms of the reference, and four free parameters 
\(\bm{b},\bm{c}\):
\begin{align*}
  \CenterTrace^\top
    &= \bm{r}_R^\top - \bm{b}^\top \FrameBasis^\times, \\ %[4pt]
  \TraceFT
    &= {\TraceFT}_R + \bm{c}, \\ %[4pt]
  \TraceFF
    &= {\TraceFF}_R + \bm{c}\otimes{\TraceTF}_R, \\ %[4pt]
  \TraceNF^\top
    &= {\TraceNF}_R^\top
       - (\bm{b}^\top {\TraceBT}_R)\,{\TraceTF}_R^\top
       + \bm{b}^\top {\TraceMF}_R, \\%[4pt]
  \TraceTF
    &= {\TraceTF}_R, \\%[4pt]
  \TraceBT
    &= {\TraceBT}_R, \\%[4pt]
  \TraceMF
    &= {\TraceMF}_R.
\end{align*}
\end{proposition}
A proof is provided in \cref{sec:proof-s3}.

\subsubsection{Class III}
We now consider a final restriction which eliminates all remaining free parameters, and furnishes a unique trace matrix: the \emph{energetic} trace. 

\begin{definition}[Class III]\label{def:class-iii}
A Class I trace is a \textbf{Class III} trace if the conjugate resultants \(\TraceTotalEnergyStress\) are exactly the classical rigid-section resultants \(\snm\):
\begin{equation}
    \bar{\bm{s}} \equiv \snm
    \label{eq:def-s3}
\end{equation}
ie, the unique trace such that generated strains \(\bm{e}\) are work conjugate to the classical resultants \(\snm\).
\end{definition}

\begin{proposition}[Class III]
All Class III traces are equivalent in the Saint~Venant class, with trace matrix given by the consistency condition:
\begin{equation}
\TraceDeDpRoot = \IesanStiff^{-\top} \IesanEnergy
\label{eq:T0-formula}
\end{equation}

\end{proposition}
\begin{proof}
For a Class I trace with $\bm{e}(x) = \TraceDeDpRoot \exp(x\IesanEvolveSP)\bm{p}$, 
the energy tangent \eqref{eq:energy-conjugacy-def} becomes:
\begin{equation*}
\delta\bm{p}\cdot \left(\mathbf{1} + x\TraceEvolveEP^\top\right) \TraceDeDpRoot^\top \bm{s}(x) 
= \delta\bm{p}\cdot \left(\mathbf{1} + x\TraceEvolveEP^\top\right) \IesanEnergy\, \left(\mathbf{1} + x\TraceEvolveEP\right) \bm{p}.
\end{equation*}
Since this must hold for all $\bm{p}$, we require:
\begin{equation*}
\TraceDeDpRoot^\top \mathbf{R}(x) = \IesanEnergy\, \left(\mathbf{1} + x\TraceEvolveEP\right).
% \label{eq:conjugacy-condition}
\end{equation*}
with \(\IesanEnergy\) furnished by \eqref{eq:def-iesan-energy}. 
Expanding both sides using $\left(\mathbf{1} + x\TraceEvolveEP\right) = \mathbf{1} + x\IesanEvolveSP$ and 
matching powers of $x$:
\begin{align*}
x^0: \qquad & \TraceDeDpRoot^\top \IesanStiff = \IesanEnergy, 
\\ %\label{eq:x0-condition} \\[4pt]
x^1: \qquad & \TraceDeDpRoot^\top \mathbf{R}_1 = \IesanEnergy\,\IesanEvolveSP.
% \label{eq:x1-condition}
\end{align*}
Consistency in \(x^0\) requires:
\begin{equation*}
\TraceDeDpRoot = \IesanStiff^{-\top} \IesanEnergy
% \label{eq:T0-formula}
\end{equation*}
The linear evolution \(x^1\) yields the consistency condition:
\begin{equation}
\IesanEnergy\,\IesanStiff^{-1} \mathbf{R}_1 = \IesanEnergy\,\IesanEvolveSP
\qquad \Longrightarrow \qquad 
\IesanStiff^{-1} \mathbf{R}_1 = \IesanEvolveSP
\label{eq:consistency-condition}
\end{equation}
This condition states that the evolution of stress resultants along 
the beam matches the universal strain evolution prescribed by  $\IesanEvolveSP$.
\end{proof}
